\subsection{Canonical sections on a rational normalized cover}
\label{subsec:rational-canonical-equality}

For a rational normalized cover, Serre duality bounds the contribution
of its singularities. The argument requires only ampleness of the divisor
and supplies the rational-cover cases in the strict genus theorem.

\begin{proposition}\label{prop:rational-normalization-canonical}
Let $X$ be a smooth projective integral surface over an algebraically
closed field $k$ of characteristic $p>0$. Let $Y\subset X$ be a
smooth integral ample divisor, and let $M\ne\OO_X$ have order $p$
with $M|_Y\simeq\OO_Y$. Suppose
\[
 2g(Y)-2=pY^2,
\]
and let $\pi:N\to X$ be the normalization of the $\mu_p$-torsor
defined by a trivialization of $M^{\otimes p}$.
Assume that $N$ is birationally rational, and put $d=Y^2$.
Choose a proper birational morphism $r:S\to N$, where $S$ is
obtained from $\PP^1\times\PP^1$ by finitely many point blowups,
and put
\[
 \mathcal T=R^1r_*\OO_S,\qquad
 \epsilon=\dim_k H^0(N,\mathcal T).
\]
Then $\mathcal T$ has finite support and
\begin{equation}\label{eq:rational-normalization-canonical-cohomology}
 \begin{gathered}
 H^1(N,\OO_N)=0,\qquad
 H^2(N,\OO_N)\simeq H^0(N,\mathcal T),\\
 \chi(\OO_N)=1+\epsilon,\qquad
 \epsilon\in\{0,1\}.
 \end{gathered}
\end{equation}
If $\epsilon=1$, then
\[
 \omega_N\simeq\OO_N,\qquad
 \bigl(\omega_X(-(p-1)Y)\bigr)^{\otimes p}\simeq\OO_X,\qquad
 K_X\equiv(p-1)Y.
\]
In particular:
\begin{enumerate}[label=\textup{(\roman*)}]
\item No such birationally rational normalization exists for $p\geq5$.
\item If $p=2$, then $\epsilon=1$ and
\[
 \begin{gathered}
 \chi(\OO_X)=1,\qquad K_X\equiv Y,\qquad
 1\leq d=K_X^2\leq9,\\
 b_1(X)=0,\qquad b_2(X)=10-d.
 \end{gathered}
\]
\end{enumerate}
Here $b_i$ denotes the dimension of $\ell$-adic cohomology for any
prime $\ell\ne p$.
\end{proposition}

\begin{proof}
The conductor identity of
Proposition~\ref{prop:normalized-mup-conductor} makes its integer
$a$ equal to one and its nonnegative defects zero.
That proposition and
Corollary~\ref{cor:normalized-equality-torsor} show that $N$ is an
integral normal projective Gorenstein surface, that $\pi$ is finite
flat and radicial of rank $p$, and that
\[
 D=\pi^*Y
\]
is a smooth integral ample Cartier divisor.
They also give an isomorphism of invertible sheaves
\[
 \omega_N\simeq
 \pi^*\omega_X\otimes\OO_N(-(p-1)D)
 \simeq\pi^*\bigl(\omega_X(-(p-1)Y)\bigr).
\]
This is the absolute dualizing sheaf of $N$.
The formula follows from finite duality and the conductor ideal,
as proved there; it is not a trace calculation in the inseparable
function-field extension.

Adjunction on $Y$ gives $K_X\cdot Y=(p-1)d$, and the projection
formula therefore gives
\[
 K_N\cdot D
 =p\bigl(K_X\cdot Y-(p-1)Y^2\bigr)=0.
\]
Both divisors in this intersection are Cartier.

We include the cohomology calculation for the stated ample
hypotheses. Since $N$ is birationally rational, its function field
is that of $\PP^1\times\PP^1$. Resolve the resulting rational map
to $X$ by point blowups \cite[Tag~0C5H]{Stacks}.
The resulting surface $S$ is smooth and projective.
The morphism $S\to X$ extends to $r:S\to N$: sections of the
normalization algebra are integral over the local rings of $S$,
and belong to them by normality. This is the extension argument
in Proposition~\ref{prop:distinguished-ruled-cover}.
The resulting map $r$ is proper and birational, and $\pi r$ is a
purely inseparable morphism of generic degree $p$.

For a point blowup of a smooth surface, the pushforward of the
structure sheaf is the structure sheaf and the positive higher
direct images vanish \cite[Lemma~54.3.4, Tag~0AGP]{Stacks}.
The two affine charts give this vanishing directly: on their
intersection, each Laurent monomial in the chart parameter belongs
to one of the two chart rings, so the degree-one \v{C}ech
differential is surjective. The analogous calculation for
$\PP^1_R$ and Leray consequently give
\[
 H^0(S,\OO_S)=k,\qquad H^1(S,\OO_S)=H^2(S,\OO_S)=0.
\]

Normality gives $r_*\OO_S=\OO_N$. Indeed proper direct-image
coherence \cite[III$_1$, Theorem~3.2.1]{EGA} makes the pushforward
a finite algebra on every affine open of $N$.
Its sections lie in the common function field and are integral
over the normal affine ring, so they belong to that ring.
Moreover $r$ is an isomorphism near each codimension-one point of
$N$ \cite[Tag~0BFP]{Stacks}.
Thus the coherent sheaf $\mathcal T$ has finite support.

For completeness, $R^ir_*\OO_S=0$ for $i\geq2$ can be checked
without assumptions on smoothness of the fibres.
The closed graph gives a projective embedding of $S$ over $N$,
and every fibre has dimension at most one.
At a closed point of $N$, choose a relative hyperplane containing
no component of its fibre, then another avoiding the finite
intersection of the first with that fibre.
The infinite field $k$ permits these choices.
Properness gives an affine neighbourhood on which the two
hyperplanes have empty common intersection with $S$.
Their complements and their intersection are affine.
The resulting two-open \v{C}ech complex has no cohomology in
degree at least two. These neighbourhoods give the asserted
higher-direct-image vanishing.

The low-degree Leray sequence now reads
\[
 \begin{split}
 0\longrightarrow H^1(N,\OO_N)
 &\longrightarrow H^1(S,\OO_S)
 \longrightarrow H^0(N,\mathcal T)\\
 &\longrightarrow H^2(N,\OO_N)
 \longrightarrow H^2(S,\OO_S).
 \end{split}
\]
The two displayed positive cohomology groups on $S$ vanish.
Since $N$ is proper and integral over $k$, one also has
$H^0(N,\OO_N)=k$. This proves the cohomology and Euler
characteristic assertions in
\eqref{eq:rational-normalization-canonical-cohomology}.

The Gorenstein surface $N$ is Cohen--Macaulay of pure dimension
two. Serre duality for such a proper scheme
\cite[Lemma~48.27.5, Tag~0FVZ, and Remark~48.27.6, Tag~0FW0]{Stacks}
gives
\[
 H^0(N,\omega_N)\simeq H^2(N,\OO_N)^\vee.
\]
If $\epsilon>0$, choose a nonzero section of $\omega_N$.
As $N$ is integral, its zero scheme is an effective Cartier
divisor unless empty. Every nonzero effective divisor has
positive intersection with the ample divisor $D$.
The equality $K_N\cdot D=0$ therefore makes the zero divisor empty.
There cannot be an isolated zero: a nonzero nonunit in a
noetherian normal local domain has a height-one zero by the
principal ideal theorem.
The section is consequently nowhere vanishing, so
$\omega_N\simeq\OO_N$. It follows that
$h^0(N,\omega_N)=h^0(N,\OO_N)=1$, and duality gives $\epsilon=1$.
This proves $\epsilon\in\{0,1\}$.

Suppose $\epsilon=1$, and put $L=\omega_X(-(p-1)Y)$.
The canonical formula gives $\pi^*L\simeq\OO_N$.
The sheaf $\mathcal E=\pi_*\OO_N$ is locally free of rank $p$.
The finite projection formula yields
\[
 \mathcal E\otimes L
 \simeq\pi_*(\pi^*L)
 \simeq\mathcal E.
\]
Taking determinants and cancelling the invertible sheaf
$\det\mathcal E$ gives $L^{\otimes p}\simeq\OO_X$.
This is the scalar-line case of the determinant norm
\cite[Tag~0BD2]{Stacks}, with no separability requirement.
Its degree on every integral curve is zero, whence
$K_X\equiv(p-1)Y$ and $K_X^2=(p-1)^2d$.

Corollary~\ref{cor:normalized-equality-euler} gives
\begin{equation}\label{eq:rational-normalization-canonical-euler}
 1+\epsilon
 =p\chi(\OO_X)-\frac{p(p-1)(p-2)}{12}d.
\end{equation}
For a prime $p\geq5$, the equality $2g(Y)-2=pd$ makes $d=2e$
even. The integer $(p-1)(p-2)$ is divisible by six.
Thus the correction term in
\eqref{eq:rational-normalization-canonical-euler} is
$p(p-1)(p-2)e/6$, an integer divisible by $p$.
It follows that $1+\epsilon\equiv0\pmod p$, contradicting
$\epsilon\in\{0,1\}$.
This proves \textup{(i)}.

For $p=2$, the correction term vanishes, so
$1+\epsilon=2\chi(\OO_X)$.
Hence $\epsilon=1$ and $\chi(\OO_X)=1$, and the preceding
determinant argument gives $K_X\equiv Y$.
The purely inseparable morphism $\pi r:S\to X$ has smooth
projective source and target. Ekedahl's comparison
\cite[Lemma~0.1.5, p.~103]{EkedahlCanonical} gives
$b_1(X)=b_1(S)=0$; finiteness of that morphism is not required
for this assertion.
The class of $Y$ in $H^2_{\mathrm{et}}(X,\mathbf Q_\ell(1))$
is nonzero, since its square has degree $d>0$.
Therefore $b_2(X)\geq1$.
Poincaré duality and Gauss--Bonnet on the smooth surface $X$
\cite[Section~45.7, especially Lemma~45.7.8, Tag~0FH0]{Stacks}
give $c_2(X)=2+b_2(X)\geq3$.
Noether's formula
\cite[Section~42.66, Tag~02UO]{Stacks} then gives
\[
 12=12\chi(\OO_X)=K_X^2+c_2(X)=d+2+b_2(X).
\]
Consequently $1\leq d\leq9$ and $b_2(X)=10-d$, as asserted
in \textup{(ii)}.
\end{proof}

\begin{corollary}\label{cor:characteristic-three-rational-residual}
In the setting of Proposition~\ref{prop:rational-normalization-canonical},
suppose $p=3$.
Then either
\[
 \epsilon=0,\qquad d\equiv4\pmod6,
\]
or
\[
 \begin{gathered}
 \epsilon=1,\qquad d=2,\qquad
 K_X\equiv2Y,\qquad \chi(\OO_X)=1,\\
 K_X^2=8,\qquad c_2(X)=4,\qquad b_2(X)=2.
 \end{gathered}
\]
In the first alternative, every singular point of $N$ is a rational
double point.
\end{corollary}

\begin{proof}
The smallest equality makes $d=2e$ even, and
\eqref{eq:rational-normalization-canonical-euler} becomes
\[
 1+\epsilon=3\chi(\OO_X)-e.
\]
Since $\epsilon\in\{0,1\}$, the case $\epsilon=0$ requires
$e\equiv2\pmod3$, or $d\equiv4\pmod6$.

In this case $\mathcal T$ has finite support and
$H^0(N,\mathcal T)=0$, so $\mathcal T=0$.
Indeed its annihilator defines a zero-dimensional finite-type
$k$-scheme, hence a finite $k$-scheme, and pushforward of its
module to $\Spec k$ is faithful.
Localizing at a singular closed point of $N$ gives a smooth
modification with $H^1(\OO)=0$.
Any other normal modification is dominated by point blowups of
that model, by Tag~0C5H. Those blowups preserve the vanishing,
and the low-degree Leray sequence injects the $H^1(\OO)$ of
the normal modification into that of the dominating smooth model.
Thus $N$ has rational singularities in the all-modifications
sense \cite[Tag~0B4N]{Stacks}.
Its singular local rings are two-dimensional normal Gorenstein
algebraic local rings. Excellence makes them Nagata with normal
completions \cite[IV, 7.8.3(ii), (iii), (v), (vi)]{EGA}.
The rational-Gorenstein criterion therefore makes these
singularities rational double points in every characteristic
\cite[Sections~54.9 and~54.12, Tags~0B4V and~0BGB]{Stacks}.

Suppose $\epsilon=1$.
The canonical and determinant conclusions give
$K_X\equiv2Y$ and $K_X^2=4d$, while the Euler formula gives
$\chi(\OO_X)=(d+4)/6$.
The smooth rational model mapping purely inseparably to $X$
again gives $b_1(X)=0$, and the ample class gives $b_2(X)\geq1$.
Noether's formula and Gauss--Bonnet therefore yield
\[
 3\leq c_2(X)
 =12\chi(\OO_X)-K_X^2
 =2(d+4)-4d
 =8-2d.
\]
The positive even integer $d$ must be two.
Substitution gives all the displayed invariants.

\end{proof}

\begin{remark}\label{rem:rational-canonical-scope}
These conclusions concern merely ample divisors and assume that the
normalization is birationally rational. They do not exclude a ruling
with positive-genus base. In characteristic two they give a necessary
canonical numerical configuration; in characteristic three they give
the two alternatives above. The strict genus theorem below excludes
both rational and positive-genus ruled possibilities when $Y$ is very
ample. It does not settle strict inequalities or higher conductor values.
When $\epsilon=1$, the nonzero higher-direct-image defect witnesses a
nonrational singularity; no rational-singularity hypothesis is imposed.
\end{remark}
